\section{收复、设省与海岛：边疆被重新安排}

一八七〇年代末，清廷同时背着几条很不相同、却彼此牵连的路：灾区的粮路从仍能通行的河道和商号向外伸展；左宗棠部的军需要越过甘肃与河西；伊犁的照会在北京、圣彼得堡和河谷之间往返；福建与台湾的港口又必须计算煤炭、船只和讯息抵达的先后。它们不能合成一部从北京向边缘展开的行政史。每一条路都要靠不同的人、货物和中介维系，也都把风险留在沿线。

战争的胜负因而只是重新安排的开端。军队进入绿洲，并不立即解决水利、市场与赋役；建省写入成案，并不使山地、港口和平原在同一日受同一种治理；条约的中外文本也不能代替关卡、牧地与市镇中实际发生的往来。晚清的边疆政治首先是一种昂贵的连接工程：朝廷试图把军令、银粮、信息和行政名义接起来，地方的运输者、商人、居民与社区则在征购、借贷、迁徙、贸易和避险中承受其不均衡的后果。

\subsection{饥荒中的粮路：救济不是一张覆盖全国的网}

\eventanchor{EQ-1876-NORTH-CHINA-FAMINE}{华北、西北饥荒（1876）}发生时，旱情并不是唯一的压力。粮价、仓储、道路和地方可动用的信用共同决定一袋粮能否离开产地、经过关卡或市镇、到达一个正在断炊的村庄。对离乡者而言，赈粮只是几个有限选择中的一个；投亲、借贷、出售物品、结伴迁徙或留在原处等待，都会随家庭劳力、可走道路和地方关系而改变。把灾区写成静待命令抵达的平面，便看不见这些被迫作出的判断。

到一八七七年，\eventanchor{EQ-1877-FAMINE-RELIEF-NETWORKS}{灾区赈济网络}把清廷的赈务、地方绅商的募捐、商路的转运和传教团体的救助接入同一场危机。它们有时在同一市镇相遇，有时只抵达各自能够取得消息、粮源和保护的地点。募捐簿上的一笔银两、奏报中的一批粮食，首先记录的是某项筹措或交付；它不能自动说明谁领到粮、一个家庭能维持多久，或邻近村庄是否被遗漏。灾民留下的自述本已稀少，迁徙者、儿童和无力立券记账者尤其难以在保存下来的文书中发声。

救济还改变了信息的价值。能够向县城、会馆或通商口岸传出消息的地点，更容易进入募捐者、报刊和外来救助者的视野；远离干道的村落未必因此没有互助，却较少留下可被后人追踪的记录。地方绅士的名册会突出捐输者，传教网络的报告会说明本机构的行动，督抚的奏报则要让上级看见可处理的灾情。把这些材料并置，不是要从中凑出一个无误差的死亡数，而是要辨认粮食、信用和消息曾在哪里相互接上，又在哪里中断。

这场饥荒与随后西征并非同一条因果线，却在同一个财政世界中争夺可用的粮、银与运输能力。省级官员、商人和地方组织既要应付本地的救济与治安，也要面对跨省筹饷和采购带来的价格、劳力与道路压力。清廷能够调动资源，并不意味着成本在各地平均分摊；正是在这种不均衡中，军事收复和灾后生计同时变得可能，也同时留下新的债务与脆弱性。

\subsection{河西饷道、绿洲接收与伊犁谈判}

\begin{event}{\eventanchor{EQ-1878-Zuo-ZONGTANG-XINJIANG}{收复新疆大部：饷道如何决定军队能够走多远}}{光绪四年（一八七八年）}{左宗棠部、清廷与新疆各地行动者 · 甘肃、河西走廊、吐鲁番、喀什噶尔及北路牧地}{清军收复大部与阿古柏政权崩解可核；战事日期、经费、人员与各地接收的范围仍须逐件复核}
左宗棠部继续西进时，前线的胜负始终系在后方能否把银粮、牲畜、道路修缮和递送文书连成一条不断裂的线。甘肃与河西不是地图上通往新疆的一段空白：地方官要筹解，商人和运输者要估量货价、牲力与治安，沿线居民则要面对征购、劳役和驻军需求。军队能抵达某座城，说明这条链在某一时刻仍可维持；它不表示每一笔饷项都如数到位，更不表示沿路社会可以不受代价地提供它。

阿古柏政权崩解后，吐鲁番、喀什噶尔、伊犁与北路牧地也没有进入同一个整齐的“战后”。军队到达、城镇市场重开、州县与驻防设置、水利和土地关系如何继续，各有自己的节奏。清廷由此重新获得安排行政的机会，却必须先决定何处优先供给、怎样与既有的地方关系打交道、哪些道路足以维持驻军。设省不是凯旋之后附加的一枚印章，而是要把临时军需转成日常征解、司法和交通时才显露出来的难题。

清方军报、军机处文书和地方呈报可以勾勒调兵、筹饷与战况的次序；它们不能替绿洲与牧地中的人说明服役、贸易或迁居如何改变。俄文材料、维吾尔文和蒙古文文书以及地方档案保留了另一组观察位置，但保存不全、地名与纪年转换亦有难处。因而，“收复”在这里应当先读作清军所确认的军事过程，而不是把不同地点此后的服从、负担和安全感预先写成结论。
\end{event}

设省后的新行政也要在旧有的生活空间中落脚。州县、驻防、屯垦和军台可以在文书上各有职责，实际运作却仍要依赖掌握水源、土地、市场和道路信息的人。一个税名的设立、一次驻军的调动或一段军台的整修，可能为某处带来新的交易机会，也可能增加征发与纠纷；不能把行政设置直接等同于居民已经获得相同的司法、教育或赋役经验。新省所要处理的，正是把战时的远距离供给改造成可以反复征解、传递与裁断的日常程序，而这种转换从来不在地图上同时完成。

军事压力转入外交，仍没有让这条饷道失去作用。清俄围绕伊犁归还和边界条件持续交涉的\eventanchor{EQ-1880-ILI-NEGOTIATIONS}{一八八〇年谈判}，发生在北京与圣彼得堡的照会网络中，也牵动伊犁河谷的牧地、城镇和跨境商路。清廷需要把西征取得的军政位置转换为可被承认的条件；俄国方面则将边境安全、商业与帝国利益写入自己的交涉语言。河谷中的商队、迁徙家庭和地方中介却未必依照两座首都的节奏行动：一项条文若要成为关卡的惯例，还要经过翻译、传达、征税和实际通行。

\eventanchor{EQ-1881-SAINT-PETERSBURG-TREATY}{《中俄伊犁条约》（1881）}使伊犁大部归还、赔款与贸易条件获得了条约形式。签署、批准和中俄文本能够确定双方作出了哪些承诺；它们并不能画出居民已经接受的精确边线，也不能把签约日变成地方接收完成日。对清廷而言，条约减轻了继续以战场方式处理伊犁的压力，却把跨境贸易、赔款和边界管理一并带进财政与行政的日常账目。对河谷社会而言，哪些货物能够过关、谁要承担税役、牧地和亲属网络如何被新手续切分，仍须从地方与多语材料中逐处追问。

条约中的赔款与通商条件也提醒人们，外交并非发生在军需之外。赔付需要筹措和解送，贸易条款需要在口岸、关卡和账册中被解释，边界管理则要不断处理来往的人、牲畜和货物。清廷和俄国的照会能清楚呈现各自要坚持的语言，却较难记录一笔交易因新手续拖延多久、一户迁徙家庭怎样重排亲属关系。这里的证据限制不只是史料学的附注；它改变了对“定界”一词的读法，使边界成为持续被执行、规避和协商的过程。

\subsection{从北圻到台湾：海防是一组港口、道路和账目的选择}

\begin{event}{\eventanchor{EQ-1883-SINO-FRENCH-WAR-BEGIN}{北圻交涉变成战争}}{光绪九年（一八八三年）}{清、法国、越南及两广边路行动者 · 北圻、红河、两广与北京}{冲突升级为战争可核；开战责任、地方控制与兵力规模不能由任一方档案裁定}
北圻的冲突升级时，红河水路、城镇、两广边路与海防港口同时决定军需和消息能够怎样移动。清廷面对的并不只是关于宗藩关系的外交措辞：边将要判断何处可驻、何处可运粮，越南朝廷和地方行动者也在法国军政推进与本地关系之间作出选择。法国拥有海军优势，并不等于它无需把炮火、河口、登陆、补给和城镇维持接在一起。各方都在有限讯息下行动，战争因而不是两份声明的自动结果。

北圻的战局也把福建、广东和台湾卷入同一场资源竞争。两广边路所需的转运与驻军，港口防务所需的煤炭、修理和船员，都会占用本可用于别处的银钱与注意力。北京发出的调兵令能够显示它试图如何排列优先次序，却不能替每一段边路说明翻译、商路、雨季和地方避险怎样改变了命令的落点。把越南只写成清法竞争的舞台，或把沿海社区只写成海军的背景，都会抹去这些相互牵制的时间。
\end{event}

\eventanchor{EQ-1884-FUZHOU-NAVAL-BATTLE}{马江海战（1884）}使福州船政、港湾防御和法国舰队在短时间内正面相遇。船只并非脱离训练、指挥、煤炭、港口和情报而独立发挥作用；一支舰队即使在纸面上拥有船只，也要能让人员、弹药、修理和讯息同时抵达。战后，海防经费的去向不再只是海图上的问题：船坞附近的工匠、运输煤粮的船户、依港口市场谋生者以及承担加派的腹地，都要在战争留下的支出中寻找位置。中法军报、船政记录和地方材料对于舰只损失与伤亡并不总能相互吻合，尤其不能用一份战报的数目概括港湾内外的全部经验。

台湾因此不是这场战争的附带地名。封锁与基隆、沪尾一带的防御迫使决策者把港口、煤源、粮运、电报与炮台放在同一张预算中计算。战后，\eventanchor{EQ-1885-TAIWAN-PROVINCE}{台湾建省（1885）}把军防、行政和工程纳入新的制度计划。刘铭传推动铁路、电报、矿务与抚垦时，计划、拨款、征地、施工和日后使用是彼此不同的步骤：淡水、基隆和台南的港市，平原移民聚落和内山各社群与这些步骤的关系不会相同。省名可以表明朝廷重新排列行政与防务的意图，不能压平岛内因土地、劳役、道路和武力接触而出现的差异。

工程尤其使行政的尺度变得可见。一段铁路或电报线要先取得资金、材料、测量和施工的条件，才可能改变何处接到消息、何处运输较快；矿务的设想也要经过土地、劳力和市场才会产生可见后果。地方档案与契约能够让我们在一些地点看见征地或劳役的痕迹，不能把这些片段推广成全岛统一的经验。内山社群的声音在留存材料中常被军政与开发文书遮蔽，正是“经营台湾”不能被写成一项没有摩擦的行政成就的原因。

\subsection{朝鲜危机、甲午战场与海峡的后果}

\begin{event}{\eventanchor{EQ-1894-FIRST-SINO-JAPANESE-WAR}{甲午开战：朝鲜、黄海与辽东的补给压力}}{光绪二十年（一八九四年）}{清、日本与朝鲜行动者 · 朝鲜半岛、黄海、辽东及沿岸港口}{开战及区域扩展可核；军事规模、责任和不同社会的经历不能由一方战报判定}
朝鲜的政治行动、驻军与外交交涉先改变了危机的条件，不能把朝鲜缩成清日竞争的一块背景板。随后黄海的舰队、辽东的道路和沿海港口把煤炭、粮食、情报与伤员运输拖入不同战场。清廷与日本都试图以调兵、筹饷和文书把分散的资源接向前线；船只是否可用、港口能否补给、道路是否容得下辎重，却不由北京或东京的一道命令决定。

这场战争的直接后果是区域秩序的军事化升级，较长的代价则落在海防财政、边地居民与跨境贸易网络上。舰队兵额、战果数字和将领责任，在中日朝文及外文记录中各有用途与计算方式；它们能够帮助确认战争怎样扩展，却不能使研究者替不同地点的居民说出相同的感受。把朝鲜、黄海和辽东并置，所要看见的正是同一场开战如何在不同的补给条件中变形。
\end{event}

一八九五年的\eventref{EQ-1895-SHIMONOSEKI-TREATY}{《马关条约》}规定了割让台湾、澎湖以及赔款和通商等条件，但海峡两岸的局面没有随着条约文字停止。清方撤离、日本登陆、港口交通与地方组织之间留下了短促而充满不确定的空隙。有人试图组织抵抗，有人要保护家人、货物和村落，有人被迫随交通线或战事移动；将这些不同处境都称作一个一致的“地方反应”，反而会再次遮住岛内的差异。

\begin{event}{\eventanchor{EQ-1895-TAIWAN-REPUBLIC-RESISTANCE}{台湾民主国与抵抗日军登陆}}{光绪二十一年（一八九五年）}{台湾地方行动者、清与日本 · 港口、道路、城镇与乡村}{建国名义及抵抗发生可核；组织范围、地方选择和战事过程须以中日及地方材料互校}
地方力量建立“台湾民主国”并抵抗日军登陆，表明签约后的主权转换仍须经过港口、道路、城镇与补给线。名义上的建国和实际能够动员、守住或离开哪些地点，并不是同一回事。清廷、日本方面和地方记录使用的政治语言也各不相同；它们可以标出组织、军队与行政所声称的行动，却不能使任何一方的措辞代表全岛居民的判断。

这一段抵抗不是条约自动产生的反应。铁路、电报和港口设施曾被设想为加强岛内联系的工具，战时却同样成为移动、传讯与争夺的条件；它们的作用取决于何处修成、何处可用以及谁能接近。割让后的创伤因此不能从一纸文本读尽，也不能只由登陆战决定。土地契约、地方档案、工程记录与日文军政材料能够让研究者在若干地点追踪征地、劳役、航线和武力接触，但原住民社群与普通居民的自述保存尤其有限，留在材料之外的经验必须被明确保留。
\end{event}

\subsection{看见连接，也看见断裂}

晚清的边疆与海岛政策并非没有效力：军队重新进入新疆，伊犁谈判形成条约，台湾的行政、通信和工程获得新的组织方式，港口与边路也被纳入更密集的军事筹算。这些安排解决了清廷在特定时刻的通行、供给与防务问题；它们的副作用同样具体。长距离饷运把征购和信用压力推向沿线，设省把土地、劳役与行政重新分配，条约和战争又使边境市镇、商队和港口社会必须适应新的手续与风险。

可以保存下来的文书多半来自正在调度、请款、谈判或归责的机构。它们善于留下命令怎样起草、银粮怎样呈报、条约怎样措辞，也较少留下未被登记的迁徙、拒绝、协商和家庭离散。地方档案、契约、价格记录、港口材料以及俄文、维吾尔文、蒙古文、越南文、日文和其他外文材料能够补足若干地点，却不构成一张没有空白的全景。正因如此，收复、建省、海战和割让都应被读成仍在进行的区域过程：国家的线路向外延伸，地方社会却从未只是线路末端的被动注脚。
